\subsection{Sprint 2}

Após a primeira entrega de código ao cliente, percebemos que era necessário fazer algumas mudanças significativas. Não conseguimos entregar nenhum dos entregáveis da sprint anterior e sabíamos que precisávamos mudar nossa abordagem e nos concentrar mais para que o projeto avançasse. Na primeira semana da nova sprint, os dois \ac{ages} IV do projeto decidiram alterar a forma como trabalharíamos. Eles dividiram as \ac{us} em pequenas tarefas e atribuíram uma tarefa a cada membro da equipe, para facilitar o acompanhamento e a assistência.

Eu me ofereci para assumir a tarefa de coletar o tempo de tela do usuário. Sabia que seria desafiador, mas isso também me motivou, pois se conseguisse desenvolvê-la, muitas funcionalidades do aplicativo estariam bem encaminhadas. A solução que encontrei foi utilizar a biblioteca \texttt{screen\_state} do Flutter, que monitora a tela do celular para verificar se está bloqueada ou desbloqueada. Além disso, utilizei o plugin \texttt{Pausable\_Timer} para controlar o tempo que o usuário passou com a tela desbloqueada, pausando o \textit{timer} quando a tela era bloqueada.

No entanto, ainda precisava enviar o tempo coletado para o banco de dados. Como o banco de dados estava sendo desenvolvido apenas nessa mesma sprint, tive pouco tempo para aprender mais sobre o Flutter, pesquisar sobre minha tarefa e entender os conceitos básicos de como nosso Firebase estava funcionando.

Durante o último final de semana da sprint, toda a equipe se concentrou em concluir suas tarefas e entregar o máximo possível para gerar um \ac{apk} a ser entregue ao cliente. Pedi a um dos \ac{ages} II do projeto que me ajudasse a coletar o tempo e enviá-lo para o banco de dados, e também para me mostrar como obter os dados referentes ao tempo de notificação configurado pelo usuário. Depois de algumas horas tentando, percebemos que não seria possível concluir essa parte, pois isso afetaria muito as tarefas dos outros membros da equipe, exigindo retrabalho, e não havia tempo suficiente para isso.

No final, conseguimos gerar o \ac{apk} e entregar tudo o que estava pendente da Sprint 1 e o que foi proposto na Sprint 2, exceto minha tarefa. Isso me deixou um pouco frustrado, mas sabia que, por ser algo mais complexo, os \ac{ages} III e IV entenderam e me ajudariam na próxima sprint para que pudesse concluir essa tarefa.